\documentclass{article}

\usepackage{amsmath,amssymb}
\usepackage{xurl}
\usepackage[hidelinks]{hyperref}

% Shared commands for notes in docs/paper-gaps/.

\DeclareUnicodeCharacter{2080}{\ifmmode{}_{0}\else\textsubscript{0}\fi}
\DeclareUnicodeCharacter{2081}{\ifmmode{}_{1}\else\textsubscript{1}\fi}
\DeclareUnicodeCharacter{2082}{\ifmmode{}_{2}\else\textsubscript{2}\fi}
\DeclareUnicodeCharacter{2099}{\ifmmode{}_{n}\else\textsubscript{n}\fi}

\newcommand{\Lean}{\textsc{Lean}}
\ExplSyntaxOn
\NewDocumentCommand{\leanid}{m}{
  \ifmmode
    \textsf{\detokenize{#1}}
  \else
    \group_begin:
    \urlstyle{sf}
    \tl_set:Nn \l_tmpa_tl {#1}
    \tl_replace_all:Nnn \l_tmpa_tl {\_} {_}
    \exp_args:NV \path \l_tmpa_tl
    \group_end:
  \fi
}
\ExplSyntaxOff
\providecommand{\tr}{\operatorname{tr}}
\newcommand{\tnleanIssueBase}{https://github.com/LionSR/TNLean/issues/}
\newcommand{\tnleanPrBase}{https://github.com/LionSR/TNLean/pull/}
\newcommand{\tnleanDocBase}{https://sirui-lu.com/TNLean/docs/}
\newcommand{\ghissue}[1]{\href{\tnleanIssueBase#1}{GitHub issue \##1}}
\newcommand{\ghpr}[1]{\href{\tnleanPrBase#1}{PR \##1}}
\newcommand{\doclink}[1]{\href{\tnleanDocBase#1.html}{\path{#1}}}

% Dirac notation and the identity operator, as in blueprint/src/macros/common.tex.
% \providecommand keeps note-local definitions authoritative.
\usepackage{dsfont}
\providecommand{\ket}[1]{|#1\rangle}
\providecommand{\bra}[1]{\langle #1|}
\providecommand{\braket}[2]{\langle #1 | #2 \rangle}
\providecommand{\ketbra}[2]{|#1\rangle\!\langle #2|}
\providecommand{\Id}{\mathds{1}}
\providecommand{\one}{\mathds{1}}
\providecommand{\C}{\mathbb{C}}
\providecommand{\MN}[1]{M_{#1}(\C)}
\providecommand{\MD}{M_D(\C)}

% External referencing for paper sources.  The build helper writes sanitized
% auxiliary files under build/paper-aux/ and excludes duplicated source labels.
\usepackage{xr}

% Import a generated paper auxiliary file with a caller-supplied label prefix.
% The candidate paths support builds from docs/paper-gaps/, the repository root,
% and build/paper-aux/ itself.
\newcommand{\importpaperaux}[2]{%
  \IfFileExists{../../build/paper-aux/#2.aux}{%
    \externaldocument[#1]{../../build/paper-aux/#2}%
  }{%
    \IfFileExists{build/paper-aux/#2.aux}{%
      \externaldocument[#1]{build/paper-aux/#2}%
    }{%
      \IfFileExists{#2.aux}{%
        \externaldocument[#1]{#2}%
      }{}%
    }%
  }%
}

% General external reference macros
\newcommand{\paperref}[2]{\ref{#1-#2}}
\newcommand{\papereqref}[2]{(\ref{#1-#2})}

% CPSV16 source (1606.00608 / MPDO-22-12-17-2.tex) external document and shortcuts
\importpaperaux{cpsv16-}{1606.00608/MPDO-22-12-17-2-tnlean}
\newcommand{\cpsvref}[1]{\paperref{cpsv16}{#1}}
\newcommand{\cpsveqref}[1]{\papereqref{cpsv16}{#1}}

% Verdict marker: \gapnote{<kind>}{<status>}. Kind is the policy
% classification (clarification, local-correction, scope-restriction,
% unfaithful, false-source, open-gap); status is open, wip (elimination
% underway), resolved, or historical. Machine-read by the site index;
% prints a verdict line.
\newcommand{\gapnote}[2]{\par\noindent\textbf{Verdict:} #1 (#2).\par}


\title{Wolf Proposition~1.6: Complete Positivity from Positivity, Commutative Side}
\author{}
\date{2026-08-05}

\begin{document}
\maketitle

\section{The source statement}

Wolf's Proposition~1.6 reads: let $T:\mathcal A\to\mathcal B$ be a positive
linear map between unital $C^*$-algebras; if either $\mathcal A$ or
$\mathcal B$ is commutative then $T$ is completely positive
\cite[Proposition~1.6]{Wolf2012Quantum}. The local source is
\path{Notes/WolfNoteTexSource/ch01_deconstructing_quantum.tex},
lines~600--612. Wolf proves only the finite-dimensional case, reduces to
$\mathcal B$ commutative (the $\mathcal A$-commutative case follows because
$T$ is completely positive iff its dual $T^*$ is), and then argues by
simultaneously diagonalizing the pairwise-commuting family
$\{T(a_{ij})\}$. His remark immediately after the proof states that the
algebra structure of $\mathcal A$ was never used, so the same argument gives
complete positivity for any positive map out of an operator system into a
commutative $C^*$-algebra.

\section{Formalized component}

\Lean's ambient objects for this proposition are linear endomorphisms of
$M_D(\mathbb C)$: \leanid{IsPositiveMap} and \leanid{IsCPMap} in
\leanid{TNLean/Channel/Basic.lean} both fix $\mathcal A=\mathcal B=M_D(\mathbb
C)$. Reading Wolf's ``$\mathcal B$ commutative'' hypothesis literally in this
ambient setting would force $M_D(\mathbb C)$ itself to be commutative, i.e.\
$D\le 1$, which is not the statement Wolf's proof establishes and not a
useful specialization. The formalization instead uses the range-only
hypothesis that Wolf's own remark licenses:
\leanid{PositiveOnAbelian.HasCommutingRange T}, defined as
$\forall\,X\,Y,\ \operatorname{Commute}(T(X))(T(Y))$, in
\leanid{TNLean/Channel/Schwarz/PositiveOnAbelian/CompletePositivity.lean}.
Since a family of pairwise-commuting elements generates a commutative
subalgebra, \leanid{HasCommutingRange T} is exactly the statement that $T$
maps into some commutative subalgebra of $M_D(\mathbb C)$ -- the hypothesis
Wolf's proof actually consumes, not an invented strengthening or weakening
of it.

\leanid{PositiveOnAbelian.isCPMap_of_isPositiveMap_of_hasCommutingRange}
proves the commutative-codomain case:
\leanid{IsPositiveMap T} together with \leanid{HasCommutingRange T} implies
\leanid{IsCPMap T}. The proof follows Wolf's route: it reduces complete
positivity to nonnegativity of the block quadratic form built from the Choi
matrix, and this reduces to the already-proved simultaneous-diagonalization
lemma \leanid{PositiveOnAbelian.quadraticForm_nonneg_of_isPositiveMap_of_}
\leanid{commuting_images} (\ghissue{5468} closes the wrapping gap left open
after that lemma was proved sorry-free).

\leanid{PositiveOnAbelian.isCPMap_of_isPositiveMap_of_hasCommutingRange_adjoint}
proves the commutative-domain case by the same duality argument Wolf invokes:
$T$ is completely positive iff its trace-pairing adjoint
\leanid{Matrix.traceAdjointMap T} is, and the codomain case applies to
\leanid{Matrix.traceAdjointMap T} when its range is commutative.
\leanid{PositiveOnAbelian.isCPMap_of_isPositiveMap_of_commutingRange_or}
combines both disjuncts into the ``either side'' statement.

\section{Relationship to the source hypothesis}

\leanid{HasCommutingRange T} is implied by, but strictly weaker than, ``the
ambient codomain algebra is commutative'' when the ambient codomain is the
full matrix algebra $M_D(\mathbb C)$ with $D>1$: a map with commutative range
need not have commutative full codomain, since $M_D(\mathbb C)$ itself is
never commutative for $D>1$. This is the correct reading of Wolf's
proposition for an ambient setting where $\mathcal A$ and $\mathcal B$ are
not themselves varied as commutative or noncommutative algebras but are
fixed as $M_D(\mathbb C)$; it matches the generalization Wolf's own remark
grants (operator system to commutative algebra), not a separately restricted
theorem. No hypothesis absent from Wolf's argument is added, so the theorems
above formalize Proposition~1.6 as Wolf actually proves it, not a
corollary of it.

\bibliographystyle{alpha}
\bibliography{references}

\end{document}
